\section{五次三维簇高亏格 GW 不变量的结构（Shuai Guo）}%
\label{sec:Structure of higher-genus GW invariants of the quintic (Shuai Guo)}

\subsection{问题}%
\label{sub:The problem (Guo)}

我们考虑五次三维簇 \(X = X_5 \subset \P^4\)，并问它包含多少条亏格 \(g\)、次数 \(d\) 的曲线。在亏格零时，答案从
\[ 2875, \qquad 609250, \qquad 317206375, \qquad \ldots \]
开始。

由于 \(X\) 是 Calabi–Yau 的，模空间 \(\Mbar_{g,n}(X,\beta)\) 的虚维数为零，且弦方程给出 \(\ab<\1, \gamma_2, \ldots, \gamma_n>_{g,n,\beta} = 0\)，除非 \((g,n,\beta) = (0,3,0)\)。因此我们如下定义初级不变量
\[ N_{g,\beta} \coloneqq \ab<>_{g,0,\beta} = \int_{[\Mbar_g(X,\beta)]^{\vir}} 1 \]
并把这些数组成生成级数
\[ F_g(Q) \coloneqq \sum_{\beta \text{ effective}} N_{g,\beta}Q^{\beta}. \]

不变量 \(N_{g,\beta}\) 是有理数，而我们真正想要的整数是 BPS 不变量 \(n_{g,\beta}\)。它们由正次数部分 \(F_g^+\) 通过多重覆盖变换定义
\[ \sum_{g \geq 0} F_g^+(Q)\lambda^{2g-2} = \sum_{\substack{g \geq 0 \\ \beta > 0}} n_{g,\beta} \sum_{k \geq 1} \frac{1}{k} \ab(2 \sin \frac{k\lambda}{2})^{2g-2}Q^{k\beta}. \]
注意，只有 \(n_{g,\beta}\) 被预期为整数，而且上面的那些数正是 \(n_{0,d}\)。

\begin{table}[htpb]
    \centering
    \caption{五次三维簇的主要 BPS 不变量 \(n_{g,d}\)}
    \label{tab:quinticbps}
    \begin{tabular}{crrrrr}
        \toprule
        \(g\) & \(d=1\) & \(d=2\) & \(d=3\) & \(d=4\) & \(d=5\) \\
        \midrule
        0 & 2875 & 609250 & 317206375 & 242467530000 & 229305888887625 \\
        1 & 0 & 0 & 609250 & 3721431625 & 12129909700200 \\
        2 & 0 & 0 & 0 & 534750 & 75478987900 \\
        3 & 0 & 0 & 0 & 8625 & \(-15663750\) \\
        4 & 0 & 0 & 0 & 0 & 49250 \\
        5 & 0 & 0 & 0 & 0 & 1100 \\
        6 & 0 & 0 & 0 & 0 & 10 \\
        7 & 0 & 0 & 0 & 0 & 0 \\
        \bottomrule
    \end{tabular}
\end{table}

我们看到 \Cref{tab:quinticbps} 的每一列最终都消失。这就是 Castelnuovo 界，对五次情形它断言当 \(g > \frac{d^2+5d+10}{10}\) 时 \(n_{g,d} = 0\)，由 Liu--Ruan~\cite{LiuRuanCastelnuovo} 证明。表中各数的整性即 Gopakumar–Vafa 猜想，由 Ionel--Parker~\cite{IonelParker2018} 证明。

在亏格零时，Candelas-de la Ossa-Green-Parkes 利用镜像上的周期积分计算了 \(F_0\)，这一点我们在 \Cref{sub:Enumerative mirror symmetry in genus zero} 中讨论过。在高亏格时，Bershadsky-Cecotti-Ooguri-Vafa 给出了一个递归，由 \(h < g\) 的 \(F_h\) 产生 \(F_g\)，但相差有限多个未知常数。我们的目标是解释对数 GLSM 如何确定这一歧义。

\subsection{亏格零与 B 模型的多项式结构}%
\label{sub:Genus zero and the polynomial structure of the B-model}

回忆 \Cref{sub:Enumerative mirror symmetry in genus zero}：对五次情形我们考虑 \(I\)-函数
\[
    I(q,z) = z\sum_{d \geq 0} q^{H/z+d} \frac{\prod_{m=1}^{5d}(5H+mz)}{\prod_{m=1}^d (H+mz)^5} = I_0 z + I_1 H + I_2 \frac{H^2}{z} + I_3 \frac{H^3}{z^2},
\]
其各分量被 Picard–Fuchs 算子
\[ D^4 - 5q(5D+1)(5D+2)(5D+3)(5D+4), \qquad D = q \odv{}{q}. \]
零化。若记 \(J_k = I_k/I_0\)，并令 \(\frac{I_1}{I_0} = \log q + \tau(q)\)，使得 \(Q = e^{I_1/I_0} = qe^{\tau(q)} = q + O(q^2)\)，则亏格零镜像定理断言
\[ F_0(Q) = \frac{5}{2}\ab(J_1 J_2 - J_3). \]

就我们的目的而言，把五次公式置于一般镜像对称的图景中是有用的。设 \(\mc{X} = \ab\{(X,\omega)\}\) 是 Calabi–Yau 三维簇的族，其 Kähler 模空间为 \(\mc{M} \ni Q = e^{2\pi i \omega}\)，再设 \(\mc{X}^{\vee} = \ab\{(X^{\vee},J)\}\) 为镜像族，其复模空间为 \(\mc{M}^{\vee} \ni q\)。Picard–Fuchs 方程的解是 \(X^{\vee}\) 上的周期，它们拼成
\[ I(z) = zI_0\varphi_0 + I_1 + I_2 z^{-1} + I_3 z^{-2}, \]
它落在 \(z\varphi_0, \varphi_{\alpha}, \varphi^{\alpha}z^{-1}, \varphi^0 z^{-2}\) 的张成中。这里 \(I_1 = \sum_{\alpha} I_1^{\alpha}\varphi_{\alpha}\)，\(I_2 = \sum_{\alpha} I_{2,\alpha}\varphi^{\alpha}\)，\(I_3 = I_{3,0}\varphi^0\)，而镜像映射仍为 \(Q = e^{I_1/I_0} = qe^{\tau(q)}\)。

我们将用到五次情形复模直线上的三个特殊点：大复结构点 \(q=0\)、锥点（conifold point）\(q = 5^{-5}\)（此处某个周期消失），以及轨形（orbifold）（或 Landau–Ginzburg）点 \(q = \infty\)。镜像映射在 \(q=0\) 处有定义，其在该处的展开给出格罗莫夫–威滕不变量；而在高亏格时，我们将在另外两点处施加边界条件。

我们现在用三点函数重新表述镜像定理。取 \(\mc{M}^{\vee}\) 上的局部坐标 \(q_i\)，设 \(\ab\{\phi_i\}\) 是 A 侧相应的基，并定义
\[ C^A_{ijk}(Q) \coloneqq I_0^2 \ab<\phi_i, \phi_j, \phi_k>^X_{0,3}, \qquad C^B_{ijk}(q) \coloneqq \int_{X^{\vee}} \partial_i \partial_j \partial_k \Omega \wedge \Omega. \]

\begin{thm}[镜像定理 {\cite{Givental1998,LLY1997}}]
    我们有 \(C^A_{ijk}(Q(q)) = C^B_{ijk}(q)\)。此外 \(C^B_{ijk} \in \Delta^{-1}\Q[q]\)，故它是 \(\mc{M}^{\vee}\) 上整体定义的亚纯函数，仅沿判别式有极点。
\end{thm}

特别地，乘以 \(I_0\) 的适当幂次并施加镜像映射之后，亏格零的生成函数成为 \(q\) 的有理函数，仅沿判别式有极点。因此它由有限多个系数决定。

\begin{exm}
    对超平面基下的五次情形，
    \[ C_{111} = \frac{5}{1-5^5 q}. \]
\end{exm}

\begin{exm}
    对 \(X = X_{3,3} \subset \P^2 \times \P^2\)，Yukawa 耦合为
    \begin{align*}
        C_{111} &= \frac{3^4 q_1 (2+3^3q_1+3^3q_2)}{\Delta}, & C_{222} &= \frac{3^4 q_2(2+3^3q_1+3^3q_2)}{\Delta}, \\
        C_{112} &= \frac{(1-3^3q_1-3^3q_2)(1+2\cdot 3^3 q_1 - 3^3 q_2)}{\Delta}, \\
        C_{122} &= \frac{(1-3^3q_1-3^3q_2)(1-3^3q_1+2\cdot3^3q_2)}{\Delta},
    \end{align*}
    其中
    \[ \Delta = 1 - 3\cdot3^3(q_1+q_2) + 3\cdot3^6(q_1^2 - 7q_1q_2+q_2^2) - 3^9(q_1+q_2)^3. \]
\end{exm}

在高亏格时，我们预期存在正规化势 \(f^A_{g,i}\)，满足
\[ f^A_{g,i}(Q(q)) = f^B_{g,i}(q) \in \Delta^{-(2g-2+n)}\Q[q], \]
其中右端由图和计算。

\subsection{BCOV 费曼规则}%
\label{sub:The BCOV Feynman rule}

B 模型费曼规则与全纯反常方程源自~\cite{BCOV1994}。对五次情形，Chang--Guo--Li 在~\cite{CGL2021} 中证明了多项式性，并在~\cite{CGL2026} 中证明了费曼规则。

我们现在把 BCOV 递归作为费曼规则来讨论。它来自辛变换量子化的一个有限维模型，如 \Cref{sub:Operator formalism and geometric quantization} 所述。设 \(\mc{H} = \mc{H}^- \oplus \mc{H}^+ \subset \mc{H}_{\infty}\) 是辛的，且 \(d+1 = \dim \mc{H}^+ = \dim \mc{H}^-\)，并置
\[ R = \begin{pmatrix} A & B \\ C & D \end{pmatrix} \in \End \mc{H}, \qquad \mc{Q}(\mbf{x},\bp) = (\bp, D^{-1}\mbf{x}) - \frac{1}{2}(\bp, D^{-1}C\bp). \]
给定总势 \(F(\hbar, \mbf{x})\)，我们用高斯积分形式地定义变换后的势
\[ f(\hbar, \mbf{x}) \coloneqq (\wh{R}F)(\hbar,\mbf{x}) = \ln \int e^{\hbar^{-1}(\mc{Q}(\mbf{x},\bp') - \mbf{x}'\bp') + F(\hbar,\mbf{x}')}\d{\mbf{x}'}\d{\bp'}. \]
按稳定相展开给出一个图和。其中 \(D^{-1}C\) 提供每条边上的传播子，\(F\) 提供顶点。

对 Calabi–Yau 三维簇，传播子按指标个数分别写成 \(S^{\varphi\varphi}\)、\(S^{\varphi}\) 与 \(S\)，而顶点是正规化的 A 模型关联函数
\[ P^A_{g,n}(Q) = I_0^{2-2g}\ab<\phi_1, \ldots, \phi_n>^X_{g,n}. \]
注意，这些与后文将出现的、求解量子微分方程的算子 \(S_{\tau}(z)\) 毫无关系。

\begin{exm}
    对五次情形在亏格二且无标记时，费曼规则为
    \begin{align*}
        f_2^B ={}& P_2 + \frac{1}{2}S^{\varphi\varphi}P_{1,1}^2 + \frac{1}{2}S^{\varphi\varphi}P_{1,2} + (S^{\varphi\varphi})^2 P_{1,1} + \frac{1}{8}(S^{\varphi\varphi})^2 P_{0,4} \\
        &+ \frac{1}{8}(S^{\varphi\varphi})^3 + \frac{1}{12}(S^{\varphi\varphi})^3 + \frac{\chi}{24}S^{\varphi}P_{1,1} + \frac{1}{2}\frac{\chi}{24}S^{\varphi}S^{\varphi\varphi} \\
        &+ \frac{1}{2}\frac{\chi}{24}\ab(\frac{\chi}{24}-1)S,
    \end{align*}
    其中 \(\chi = \chi(X_5) = -200\)。这是传播子的一个三次多项式。
\end{exm}

我们可以如 \Cref{sub:R-matrix action} 那样把它解释为稳定图和。顶点 \(v\) 贡献低亏格关联函数 \(P_{g_v, \on{val}(v)}\)；内边把两个插入与一个传播子缩并；腿记录其余的插入。每个图以 \(\frac{1}{\ab|\Aut \Gamma|}\) 加权。同往常一样
\[ g(\Gamma) = h^1(\Gamma) + \sum_v g_v, \qquad 2g_v - 2 + \on{val}(v) > 0, \]
而 \(f_g^B\) 是亏格 \(g\) 的所有稳定图的贡献之和。

\subsection{全纯反常方程}%
\label{sub:The holomorphic anomaly equation}

我们现在把递归改写为微分方程。B 模型振幅 \(\mc{F}_g(z,\bar{z})\) 不是全纯的，当 \(g \geq 2\) 时其非全纯性由下式给出
\[ \partial_{\bar{i}} \mc{F}_g = \frac{1}{2}\bar{C}_{\bar{i}}^{jk}\ab( D_j D_k \mc{F}_{g-1} + \sum_{r=1}^{g-1} D_j \mc{F}_r D_k \mc{F}_{g-r} ). \]
这里 \(C_{ijk}\) 是 Yukawa 耦合，\(\bar{C}_{\bar{i}}^{jk}\) 由 \(\bar{C}_{\bar{i}\bar{j}\bar{k}}\) 用 Weil–Petersson 度量上升两个指标得到，\(D_j\) 是 Hodge 线丛相应幂次上的协变导数。右端的两项对应 \(\Mbar_g\) 的两个边界层，即一个非分离结点与一个分离结点。\(\frac{1}{2}\) 用于忘掉两支的顺序。

若选取局部 \(\bar{\partial}\)-原函数
\[ \partial_{\bar{i}}S^{jk} = \bar{C}_{\bar{i}}^{jk}, \qquad \partial_{\bar{i}}S^j = G_{i\bar{\ell}}S^{\ell j}, \qquad \partial_{\bar{i}}S = G_{i\bar{\ell}}S^{\ell}, \]
则模去一个全纯项后，反常方程成为多项式偏微分方程
\[ \pdv{\mc{F}_g}{S^{jk}} = \frac{1}{2}\ab( D_jD_k \mc{F}_{g-1} + \sum_{r=1}^{g-1} D_j \mc{F}_r D_k \mc{F}_{g-r} ). \]
这些传播子仅在相差全纯张量的意义下有定义，这一自由度正是全纯歧义的来源。

\begin{rmk}
    在我们将用于五次情形的单参数正规化变量下，上式写成
    \[ -\partial_0 P_g = \frac{1}{2}P_{g-1,2} + \frac{1}{2}\sum_{r=1}^{g-1}P_{r,1}P_{g-r,1}, \qquad -\partial_1 P_g = 0. \]
\end{rmk}

为理解图和为何满足该方程，注意在无标记的势中，非全纯生成元只出现在边因子中。因此沿传播子方向求导会挑出一条边并将其割断。两个半边成为所得低亏格顶点上的标记，由除子方程，协变导数正对应于这些标记。割开一个圈得到 \(P_{g-1,2}\)，割开一条分离边得到 \(P_{g_1,1}P_{g_2,1}\)。

\begin{rmk}
    为明确起见，这一蕴含只是单向的。BCOV 图规则蕴含反常方程，但反常方程本身并不确定图规则。为恢复同一个解，我们还需要传播子的约定、低亏格顶点数据以及边界条件。
\end{rmk}

若所有低亏格顶点都已知，则右端已知，于是我们可以对非全纯生成元积分，得到
\[ P_g = P_g^{\mr{particular}} + f_g(Y), \]
其中有限生成给出全纯歧义
\[ f_g(Y) = P_g|_{A = B_1 = B_2 = B_3 = 0} = \sum_{i=0}^{3g-3}a_{i,g}Y^i, \qquad Y = (1-5^5q)^{-1}. \]
我们看到反常方程在每个亏格留下 \(3g-2\) 个未知系数。

\begin{rmk}
    对五次情形，次数零数据连同轨形正则性消去了这些自由度中的 \(\ceil{\frac{3g-3}{5}}+1\) 个，而锥点缝隙（conifold gap）条件再给出 \(2g-2\) 个条件，于是剩下
    \[ 3g-2 - \ab(\ceil*{\frac{3g-3}{5}}+1) - (2g-2) = \floor*{\frac{2g-2}{5}} \]
    个参数，它们是真正新的有效不变量。这恰好是 \Cref{eqn:balancing condition} 所预言的 \(\infty\)-顶点未知数的个数，即满足 \(\beta \leq \frac{2g-2}{5}\) 的 \(\beta\) 的个数。
\end{rmk}

于是我们得到物理学家提出的如下猜想性算法：BCOV 费曼规则、Landau–Ginzburg 点处的轨形正则性（LG/CY 对应的一种形式，对五次情形这是 Guo-Ruan-Janda 的进行中工作），以及锥点缝隙条件。此外我们还可以加上 Gopakumar–Vafa 整性与 Castelnuovo 界。

\subsection{四种理论}%
\label{sub:Four theories}

回忆 \Cref{sec:Log GLSM (Qile Chen)}：对数 GLSM 带有两个虚圈。约化圈恢复五次格罗莫夫–威滕理论，而典范等变圈容许一个单位根特殊化，它是半单的，因而可用 Givental–Teleman 计算。二部图和
\[ \Omega^{\mr{GW}}_{g,n} = \sum_{\Gamma \text{ bipartite}} \frac{1}{\ab|\Aut \Gamma|} \ab( \prod_{v \in V_{\infty}(\Gamma)} c^{\mr{eff}}_{g_v,\beta_v}) \Cont^{\mr{can}}_{\Gamma} \]
用后者表达前者，代价是在 \(\infty\)-顶点处引入有效不变量 \(c^{\mr{eff}}_{g,\beta}\)。用形式参数 \(\mbf{C} = \ab\{c_{g,\beta}\}\) 替代这些不变量，就得到 \textit{扩张 \(\mbf{C}\)-族}，其顶点数据由亚纯与 rubber 对数 GLSM 理论提供。

我们将考虑以下四种理论。
\begin{itemize}
    \item \textit{\(t\)-扭几何理论} 置 \(\lambda_i = 0\) 并保留辅助权重 \(t\)。它同时有稳定映射版本与稳定商版本，其中商版本给出下文亏格二的局部化公式。
    \item \textit{形式五次}（formal quintic），即形式 \(\lambda\)-扭理论，置 \(t = 0\) 并令 \(\lambda_i = \zeta_5^i \lambda\)。它是半单 CohFT，故具有通用且可计算的 \(R\)-矩阵。
    \item \textit{扩张 \(\mbf{C}\)-族} 在两者之间插值，同时带有形式顶点与亚纯顶点的二部图就生活在其中。
    \item \textit{五次格罗莫夫–威滕理论} 由置 \(\mbf{C} = \mbf{C}^{\mr{eff}}\) 再令 \(\lambda \to 0\) 得到，它正是我们想要计算的理论。
\end{itemize}

\begin{warn}
    在高亏格时，\(t\)-扭稳定商理论的 \(t^0\) 系数并不等于形式 \(\lambda\)-扭理论的 \(\lambda^0\) 系数。
\end{warn}

\subsection{形式五次}%
\label{sub:The formal quintic}

形式五次及其全纯反常方程在~\cite{LhoPandharipande2019} 中被研究；与扩张理论的比较则在~\cite{GJRstructure} 中展开。

设 \(E = \mc{O}_{\P^4}(5)\)，并设 \(\pi \colon \mc{C} \to \Mbar_{g,n}(\P^4,d)\) 为泛曲线。\(t\)-扭圈为
\[ e(R\pi_* f^* E \otimes [1]) \cap [\Mbar_{g,n}(\P^4,d)]^{\vir}, \]
其中 \([1]\) 的等变第一陈类为 \(t\)。若改为从完全 \((\C^{\times})^5\)-等变圈出发，形式五次是单位根特殊化
\[ [\Mbar_{g,n}(\P^4,d)]^{\vir,\lambda} \coloneqq \ab( e(R\pi_*f^*E \otimes [1]) \cap [\Mbar_{g,n}(\P^4,d)]^{\vir}_{(\C^{\times})^5})\Big|_{t=0, \; \lambda_i = \zeta_5^i \lambda}, \]
其中 \(\zeta_5 = e^{2\pi i/5}\)。关系 \(H^5 = \lambda^5\) 使量子积一般地成为半单的，同时把 \(\lambda\) 保留为参数。

态空间、单位与配对由
\[ \mc{H}_{\lambda} = \Q(\lambda)[H]/(H^5-\lambda^5), \quad \1 = 1, \quad (\gamma_1,\gamma_2)^{\lambda} = \int_{\P^4, (\C^{\times})^5} e(\mc{O}(5))\gamma_1\gamma_2 \Big|_{\lambda_i = \zeta_5^i\lambda}, \]
给出，而 CohFT 类为
\[ \Omega^{\lambda}_{g,n}(\gamma_1, \ldots, \gamma_n) = \sum_{d \geq 0} q^d \rho_* \ab( \prod_{i=1}^n \ev_i^*(\gamma_i) \cap [\Mbar_{g,n}(\P^4,d)]^{\vir,\lambda}), \]
其中 \(\rho\) 忘掉映射并稳定化定义域。若以 \(\tau(q)H\) 作平移，其中 \(\tau(q) = I^{\lambda}_1(q)/I^{\lambda}_0(q)\)，则平移后的理论 \(\Omega^{\lambda,\tau}\) 是半单的。把系数域扩张到 \(\Q(\zeta_5,\lambda)\) 之后，幂等基满足
\[ \omega^{\lambda}_{g,n}(e_{\alpha_1}, \ldots, e_{\alpha_n}) = \sum_{\alpha}\Delta_{\alpha}^{g-1}\prod_i \delta_{\alpha,\alpha_i}. \]
于是 \Cref{thm:teleman} 给出
\[ \Omega^{\lambda,\tau}_{g,n} = \sum_{\Gamma \in G_{g,n}} \frac{1}{\ab|\Aut\Gamma|}\iota_* \Cont_{\Gamma} \]
其中装饰如 \Cref{sub:R-matrix action}。携带 \(\gamma\) 的腿取 \(R^{-1}(\bar{\psi})\gamma\)，边取
\[ V_R(z_1,z_2) = \frac{\sum_{\alpha} e_{\alpha}\otimes e^{\alpha} - \sum_{\alpha} R^{-1}(z_1)e_{\alpha} \otimes R^{-1}(z_2)e^{\alpha}}{z_1+z_2}, \]
顶点则取拓扑场论 \(\omega\)，缺失的保单位部分通过 \(T(z) = z(\1 - R^{-1}(z)\1)\) 插入到被遗忘的标记处。

\begin{rmk}
    我们看到边因子没有极点。置 \(z_2 = -z_1\) 并以形式 \(R^{-1}(z) = R^*(-z)\) 使用辛条件，分子就消失。因此它可被 \(z_1+z_2\) 整除，\(V_R\) 是两个余切线变量的幂级数。\(z_1^p z_2^q\) 的系数是置于结点处并配以 \(\bar{\psi}_1^p \bar{\psi}_2^q\) 的类，把它与相邻的两个拓扑场论张量缩并就把两支粘合起来。
\end{rmk}

\subsection{如何计算 \(R\)-矩阵}%
\label{sub:How to compute the R-matrix (Guo)}

为计算 \(R\)-矩阵，我们从形式等变 \(I\)-函数开始
\[ I^{\lambda}(q,z) = z\sum_{d\geq0}q^d \frac{\prod_{m=1}^{5d}(5H+mz)}{\prod_{m=1}^d ((H+mz)^5 - \lambda^5)} = zI_0 + I_1^{\lambda}H + z^{-1}I_2^{\lambda}H^2 + \cdots. \]
到 \(H\) 为止的系数与 \(\lambda\) 无关，故 \(I_1^{\lambda} = I_1 - I_0 \log q\)，且镜像平移 \(\tau(q) = I_1^{\lambda}(q)/I_0(q)\) 不含对数。置
\[ D = q\odv{}{q}, \qquad Y = (1-5^5q)^{-1}, \qquad L = Y^{1/5} \]
并引入基
\[ \phi_k = I_{1,1}\cdots I_{k,k}H^k, \qquad \bar{\phi}_k = I_0 L^{-k-1}\phi_k, \]
其中对角 Birkhoff 因子由如下递归定义
\[ I_{p,0} = I_p, \qquad I_{p,r} = \frac{I_{p-1,r-1}}{I_{r-1,r-1}} + D \ab( \frac{I_{p,r-1}}{I_{r-1,r-1}}), \quad 1 \leq r \leq p, \]
于是特别地 \(I_{1,1} = 1+D\tau\)。幂等元与典范坐标由
\[ e_{\alpha} = \frac{1}{5}\sum_{i=0}^4 h_{\alpha}^{-i}\frac{\phi_i}{L^i}, \qquad u_{\alpha} = h_{\alpha}\int_0^q (L(s)-1)\frac{\d{s}}{s}, \qquad h_{\alpha} = \zeta_5^{\alpha}\lambda. \]
给出。

我们将用量子微分方程确定 \(R\)-矩阵对 \(q\) 的依赖。基本解满足
\[ zDS_{\tau}(z) + S_{\tau}(z)H = M_{\tau}S_{\tau}(z), \qquad M_{\tau} = \dot{\tau}\star, \quad \dot{\tau} \coloneqq (1+D\tau)H = I_{1,1}H, \]
而在正规化典范方向 \(\alpha\) 上，稳定相渐近给出
\[ (S_{\tau})_{i\bar{\alpha}}(z)C_{\alpha}(z)^{-1} = e^{u_{\alpha}/z}R_{i\bar{\alpha}}(z), \]
于是
\[ (zD+h_{\alpha}L)R_{i\bar{\alpha}} = \sum_j (M_{\tau})_i^j R_{j\bar{\alpha}}. \]
这里 \(C_{\alpha}(z)\) 是 \(q=0\) 处的量子 Riemann–Roch（等价地，Gamma 函数）因子，它与辛条件一起确定微分方程无法察觉的那些常数。与 \Cref{sub:How to compute the R-matrix} 末尾一样，微分方程给出递归，而振荡积分给出初始条件。

记 \(R(z) = \mr{Id} + R_1 z + R_2 z^2 + \cdots\)，我们就可以用上面的方程由前面各项逐个确定 \(R_{k+1}\)。有两个特点使五次情形可以处理。第一，单位根基有五个扇区，一条 mod \(5\) 的荷规则在每一行只允许一种剩余模式。于是大多数矩阵元在任何计算开始之前就已消失。第二，每个系数都落在有限生成环
\[ \mc{R} = \Q[A,B_1,B_2,B_3,Y], \qquad A = \frac{DI_{1,1}}{I_{1,1}}, \quad B_k = \frac{D^k I_0}{I_0}, \]
中，因此稳定图和归结为有限的代数计算。

\begin{rmk}
    特别地，由维数原因，固定的亏格只需要 \(R\) 的有限多项。在 \(\Mbar_{g,n}\) 上只有总次数至多 \(3g-3+n\) 的类有贡献，其余的都积分为零，故对无标记的势我们只需 \(R_0, \ldots, R_{3g-3}\)。亏格 \(4\) 到 \(z^9\) 为止，亏格 \(5\) 到 \(z^{12}\) 为止。
\end{rmk}

